\section{Layer L3: domain ontologies and vocabularies}
\label{sec:layer3}

L3 is where meaning becomes domain-specific: what a quantity \emph{is}, what an
observation \emph{is}, what an air handling unit or a distillation column
\emph{is}. Two of these, units and observations, turn out to carry
disproportionate weight for agents.

\subsection{QUDT: the unit calculus}
\label{sec:qudt}

QUDT \cite{qudt} models quantity kinds, units, dimension vectors, systems of
units and constants. Its essential construction is that a unit carries
(i) a \emph{quantity kind} with a seven-dimensional vector over the SI base
quantities and (ii) an affine conversion to the coherent SI unit
(\texttt{qudt:conversionMultiplier}, \texttt{qudt:conversionOffset}). Two
quantities are commensurable if and only if their dimension vectors are equal;
conversion is then arithmetic.

\paragraph{Why this is the highest-value single vocabulary for industrial
agents.} Industrial plants are unit-heterogeneous as a matter of history, not of
carelessness: a retrofitted line keeps its \texttt{psi} and \texttt{degF}
instruments while the new line is metric, and both feed the same MES. The
failure modes this creates are exactly the ones type systems cannot see:

\begin{itemize}
  \item \textbf{Dimensional nonsense} (\textsc{f-dim}), writing a temperature
  into a pressure setpoint. Catchable by comparing dimension vectors, and by
  nothing weaker.
  \item \textbf{Silent scale error} (\textsc{f-unitscale}), writing $58$ where
  the tag expects \texttt{bar} but the agent was thinking \texttt{psi}. The value
  is a plausible number of the right dimension. A flat catalogue that checks the
  number against the range in whatever unit was stated accepts it. Only the
  conversion catches it.
  \item \textbf{False alarms}, the mirror image, and the one practitioners
  underestimate. A legal setpoint expressed as $250$~\texttt{kPa} for a tag
  ranged $0.5$--$4.5$~\texttt{bar} is correct; a validator without a unit
  calculus rejects it. In Experiment~E4 this is the entire source of tier~G2's
  \FPRGTwo{} false-alarm rate, and in Experiment~E5 it is why G2 forbids
  $100\% - \ActRecGTwo{}$ of legal actions.
\end{itemize}

Affine units make the point sharpest. \texttt{degC} to \texttt{degF} requires an
offset; a multiplier-only conversion table, the common shortcut, is simply
wrong for temperature, the most-used quantity in process automation. QUDT is the
only surveyed technology scoring $3$ on U3.

\paragraph{Practical note.} An agent does not need all of QUDT. The slice
required is small: for each unit in play, a symbol, a quantity kind, a dimension
vector, a multiplier and an offset. In our implementation that block costs a few
hundred tokens and is treated as \emph{schema}, reserved out of the context
budget before payload is added, without that reservation, unit-conversion
tasks fail not because the information is absent but because it is at the end of
the queue.

Alternatives exist, the Ontology of Units of Measure (OM), UCUM, UN/CEFACT
codes, but only QUDT and OM carry the dimension algebra. UN/CEFACT codes, which
\opcua{} and NGSI-LD reference, are identifiers without a calculus: necessary,
not sufficient.

\subsection{SOSA/SSN: observation semantics}

SOSA/SSN (W3C Recommendation) separates the \emph{sensor}, the
\emph{observable property}, the \emph{feature of interest}, the
\emph{observation} and the \emph{result} \cite{haller2019ssn}. That five-way separation is more
useful to an agent than it first appears, because it makes the following
distinct and separately checkable: what was measured, of what thing, by what
device, when, and with what value.

\paragraph{Agent affordance.} Three failure modes collapse without it. Answering
about the \emph{property} when asked about the \emph{asset} (or vice versa);
quoting a value without a time (\textsc{f-stale}); and asserting a number that
the observation does not support (\textsc{f-value}, eight cases in our corpus,
detectable only at tier~G4). Once an answer is modelled as an observation with a
result, ``is this answer faithful to the data?'' becomes a graph query rather
than a judgement call, and that is the basis of a genuinely automatable
faithfulness check for numeric answers, which is otherwise the hardest part of
RAG evaluation.

\subsection{Brick and SAREF}

\emph{Brick} \cite{balaji2016brick} is an RDF/OWL ontology of building
equipment, points and locations with relations such as \texttt{hasPoint},
\texttt{feeds} and \texttt{isLocatedIn}. Its design lesson generalises well beyond
buildings: it deliberately restricts itself to a small, closed set of relations
and a deep class taxonomy, which keeps queries writable by non-experts, a real
advantage when the query author is a language model. (The additive-tagging
alternative, Project Haystack \cite{haystack}, has largely converged with Brick
and adds no distinct agent affordance.)

\emph{SAREF} (ETSI TS~103~264 \cite{etsi2020saref}) plus its vertical extensions
(SAREF4BLDG, SAREF4ENER, SAREF4INMA for manufacturing, and others) models
devices, functions, commands, states and services. Its \texttt{Function}/\texttt{Command}/
\texttt{State} triad gives it a genuine, if shallow, behavioural model
(U5~$=2$), rare among L3 ontologies.

\paragraph{Assessment for industrial agents.} Brick's coverage of process
equipment is thin, and neither covers a bottling line. The transferable
contribution is structural: a small closed relation set, a deep class hierarchy,
and a convention for attaching points to equipment. Our reference ontology follows
exactly that pattern and aligns to Brick where a class exists.

\subsection{DEXPI: process connectivity}

\emph{DEXPI} \cite{dexpi} standardises P\&ID exchange: equipment, piping,
instrumentation and their connectivity, with an RDF rendering available. For an
agent doing process reasoning, ``what is between the pump and the tank?'', ``which
valves isolate this section?'', DEXPI is the only standard that carries the answer,
and its connectivity graph is directly the structure that Experiment~E2 shows to
be decisive for topology tasks. The more ambitious process ontologies
(ISO~15926 \cite{iso15926} and IOF/BFO \cite{iof}) are rigorous but costly to use
and thinly adopted, so we note them without assessing them as agent infrastructure.

\subsection{Summary of L3}

L3 contributes what neither L1 nor L2 can: the unit calculus (QUDT), observation
and provenance structure (SOSA/SSN), and domain taxonomies that let an agent
generalise from one plant to another. Its weakness is coverage, there is no
Brick for a bottling line, and that gap, rather than any technical obstacle, is
the main reason industrial agent projects still start by building a bespoke
ontology. We return to this in \S\ref{sec:challenges}.
